\documentclass[11pt]{article}

\usepackage[margin=1in]{geometry}
\usepackage{amsmath,amssymb,amsthm}
\usepackage{hyperref}

\newtheorem{theorem}{Theorem}
\newtheorem{lemma}{Lemma}
\newtheorem{corollary}{Corollary}
\theoremstyle{remark}
\newtheorem{remark}{Remark}

\newcommand{\C}{\mathbb C}
\newcommand{\R}{\mathbb R}
\newcommand{\Fix}{\operatorname{Fix}}

\title{Two Positive Subcases Of The Complex Real-Form Question\\
For Complemented Subspaces Of Complex Banach Lattices}
\author{Run \texttt{fa\_banach\_001}, agent \texttt{agent\_lane\_06}}
\date{June 16, 2026}

\begin{document}
\maketitle

\begin{abstract}
De Hevia, Martinez-Cervantes, Salguero-Alarcon, and Tradacete ask in
Remark 5.5 of \emph{A counterexample to the complemented subspace problem in
Banach lattices} whether every complemented subspace of a complex Banach
lattice is linearly isomorphic to the complexification of some real Banach
space.  We record several positive subcases and normalization criteria.  First,
every finite-codimensional complex subspace of a complex Banach lattice has
such a real form.  Second, any complex Banach space with an unconditional
Schauder basis has such a real form.  Third, an anti-linear automorphism can be
normalized into a conjugation whenever the square has a suitable holomorphic
square root.  This remains true after splitting off finitely many bad spectral
values with finite-dimensional Riesz subspaces, and hence in an important
compact-perturbation case.  In particular, the same
conclusion holds for the range of any projection which is sufficiently close
to its conjugate projection.  The full infinite-codimensional question remains
open.
\end{abstract}

\section{Source Question}

The source paper solves the complemented subspace problem for real and complex
Banach lattices in the negative.  At the end of Section 5, the authors isolate
a remaining question:

\begin{quote}
Regarding the CSP for complex Banach lattices, one question that remains open
is whether every complemented subspace of a complex Banach lattice is linearly
isomorphic to the complexification of some real Banach space.
\end{quote}

They also point out that a complexification is necessarily isomorphic to its
complex conjugate, while the known examples not isomorphic to their conjugates
fail Gordon--Lewis local unconditional structure and hence cannot be
complemented subspaces of complex Banach lattices.

\section{Definitions}

A \emph{conjugation} on a complex Banach space \(Z\) is a bounded anti-linear
operator \(J:Z\to Z\) such that \(J^2=I\).  The canonical example is the
conjugation on the complexification \(E_\C\) of a real Banach space \(E\).

We say that a complex Banach space \(Y\) has a \emph{real form} if it is
complex-linearly isomorphic to the complexification of a real Banach space.

\section{Proof Intuition}

The source question asks whether a complemented complex subspace must remember
some real structure.  A real structure is exactly a conjugation up to
isomorphism.  Thus the strategy is to manufacture a bounded anti-linear
involution on the subspace.

For finite codimension, this is possible because the defining quotient map has
finite-dimensional range.  One can choose a finite-dimensional complement
spanned by vectors from the ambient real part, so the complement is invariant
under conjugation.  Projecting the ambient conjugation back along this
complement gives a conjugation on the subspace.

For a space with an unconditional basis, the conjugation is even more concrete:
write a vector in coordinates and conjugate the coefficients.  The
unconditionality constant controls all unimodular coordinate multipliers, so
this coefficient conjugation is bounded.

The next obstruction is spectral rather than geometric.  If a space is
isomorphic to its complex conjugate, it carries an anti-linear automorphism
\(A\).  This is not automatically a conjugation, since \(A^2\) may be a
nontrivial complex-linear automorphism.  When \(A^2\) has a square root chosen
compatibly with conjugation, functional calculus renormalizes \(A\) to an
anti-linear involution.

For arbitrary codimension, a projection \(P\) onto the subspace need not
commute with the ambient conjugation \(J\).  However, if \(P\) is close to its
conjugate \(JPJ\), the anti-linear map \(PJ\) on the range has square close to
the identity.  Taking a bounded square root of this square normalizes \(PJ\)
to an actual involution.

\section{A Basic Real-Form Lemma}

\begin{lemma}\label{lem:conjugation-real-form}
Let \(Y\) be a complex Banach space.  If \(Y\) admits a bounded conjugation
\(C\), then \(Y\) is complex-linearly isomorphic to the complexification of a
real Banach space.
\end{lemma}

\begin{proof}
Let \(E=\Fix(C)=\{y\in Y:Cy=y\}\), regarded as a real Banach space with the
inherited norm.  Equip \(E_\C\) with the usual complexification norm
\[
  \|e+if\|_{E_\C}=\sup_{\theta\in\R}\|e\cos\theta+f\sin\theta\|_E .
\]
Every \(y\in Y\) has the unique decomposition
\[
  y=\frac{y+Cy}{2}+i\frac{y-Cy}{2i},
\]
with both summands inside \(E\).  Thus the natural map
\[
  E_\C\ni e+if\longmapsto e+if\in Y
\]
is a complex-linear bijection.  It is bounded since
\(\|e+if\|_Y\leq \|e\|+\|f\|\leq 2\|e+if\|_{E_\C}\).  Its inverse is bounded
because the two real components above are controlled by
\((1+\|C\|)\|y\|/2\), while
\(\|e+if\|_{E_\C}\leq \|e\|+\|f\|\).  Therefore it is an isomorphism.
\end{proof}

\section{Finite-Codimensional Subspaces}

\begin{theorem}\label{thm:finite-codim}
Let \(E\) be a real Banach space and let \(Z=E_\C\) be its complexification,
with canonical conjugation \(J\).  If \(Y\subset Z\) is a closed complex
subspace of finite codimension, then \(Y\) is complex-linearly isomorphic to
the complexification of a real Banach space.
\end{theorem}

\begin{proof}
Let \(m=\dim_\C Z/Y\).  The case \(m=0\) is immediate, so assume \(m>0\).
Choose a bounded surjective complex-linear map
\[
  F:Z\to \C^m
\]
with kernel \(Y\).  The complex span of \(F(E)\) is all of \(\C^m\), because
every element of \(Z=E+iE\) is a complex linear combination of elements of
\(E\).  Hence we may choose \(u_1,\ldots,u_m\in E\) such that
\(F(u_1),\ldots,F(u_m)\) is a complex basis of \(\C^m\).

Put \(M=\operatorname{span}_\C\{u_1,\ldots,u_m\}\).  Then \(M\cap Y=\{0\}\)
and \(Z=Y\oplus M\).  Since each \(u_j\) belongs to the real part \(E\), the
subspace \(M\) is invariant under \(J\).  Let \(R:Z\to Y\) be the bounded
projection associated with the direct sum \(Z=Y\oplus M\).

Define \(C:Y\to Y\) by
\[
  Cy=R(Jy).
\]
This map is bounded and anti-linear.  We claim that \(C^2=I_Y\).  Indeed, for
\(y\in Y\), write \(Jy=Cy+m\) with \(m\in M\).  Applying \(J\) gives
\[
  y=J(Cy)+Jm.
\]
Since \(M\) is \(J\)-invariant, \(Jm\in M\).  Applying the projection \(R\) to
the last equality yields
\[
  y=R(J(Cy))=C(Cy).
\]
Thus \(C\) is a bounded conjugation on \(Y\).  Lemma
\ref{lem:conjugation-real-form} completes the proof.
\end{proof}

\begin{corollary}
Every finite-codimensional complex subspace of a complex Banach lattice is
linearly isomorphic to the complexification of a real Banach space.
\end{corollary}

\begin{proof}
By definition, a complex Banach lattice is the complexification of a real
Banach lattice.  Apply Theorem \ref{thm:finite-codim}.
\end{proof}

\section{The Unconditional-Basis Subcase}

\begin{theorem}\label{thm:unconditional-basis}
Let \(Y\) be a complex Banach space with an unconditional Schauder basis.
Then \(Y\) is complex-linearly isomorphic to the complexification of a real
Banach space.
\end{theorem}

\begin{proof}
Let \((e_n)\) be an unconditional basis for \(Y\).  There is a constant \(K\)
such that every coordinate multiplier
\[
  M_\lambda\left(\sum_n a_ne_n\right)=\sum_n \lambda_n a_n e_n,
  \qquad |\lambda_n|\leq 1,
\]
defines a bounded complex-linear operator with norm at most \(K\).

Define
\[
  C\left(\sum_n a_ne_n\right)=\sum_n \overline{a_n}e_n .
\]
For a fixed vector \(y=\sum_n a_ne_n\), choose
\(\lambda_n=\overline{a_n}/a_n\) when \(a_n\neq 0\), and \(\lambda_n=0\)
otherwise.  Then \(Cy=M_\lambda y\), so \(\|Cy\|\leq K\|y\|\).  Thus \(C\)
is bounded.  It is anti-linear and satisfies \(C^2=I\).  Lemma
\ref{lem:conjugation-real-form} now gives the conclusion.
\end{proof}

\begin{corollary}
Any complemented subspace of a complex Banach lattice which has an
unconditional Schauder basis is linearly isomorphic to the complexification of
a real Banach space.
\end{corollary}

\begin{remark}
This includes the basis cases but not arbitrary unconditional finite-dimensional
decompositions.  Kalton's spaces \(Z_2(\alpha)\), which are not isomorphic to
their complex conjugates, show why block decompositions alone are a delicate
frontier rather than an automatic real-form criterion.
\end{remark}

\section{A Spectral Normalization Criterion}

\begin{theorem}\label{thm:spectral-normalization}
Let \(Y\) be a complex Banach space.  Suppose that \(A:Y\to Y\) is a bounded
invertible anti-linear operator.  If the spectrum of the complex-linear
operator \(L=A^2\) is contained in \(\C\setminus(-\infty,0]\), then \(Y\) is
complex-linearly isomorphic to the complexification of a real Banach space.
\end{theorem}

\begin{proof}
Since \(A\) is anti-linear and \(L=A^2\), we have \(AL=LA\).  The spectral
assumption lets us define the principal inverse square root
\[
  S=L^{-1/2}
\]
by holomorphic functional calculus on the slit plane
\(\C\setminus(-\infty,0]\).  The principal branch satisfies
\[
  f(\overline z)=\overline{f(z)}
\]
for \(f(z)=z^{-1/2}\).  By approximating \(f\) uniformly on a conjugation
symmetric neighborhood of \(\sigma(L)\) by polynomials with real coefficients,
or equivalently by the symmetry of the Cauchy integral formula, one obtains
\[
  AS=SA .
\]
Now put \(C=SA\).  Then \(C\) is bounded and anti-linear, and
\[
  C^2=SASA=S^2A^2=L^{-1}L=I .
\]
Thus \(C\) is a bounded conjugation.  Lemma
\ref{lem:conjugation-real-form} completes the proof.
\end{proof}

\begin{corollary}\label{cor:projection-spectral}
Let \(Z\) be a complex Banach space with a bounded conjugation \(J\), let
\(P:Z\to Z\) be a bounded complex-linear projection, and set \(Y=PZ\).  If
\[
  A=PJ|_Y:Y\to Y
\]
is invertible and \(\sigma(A^2)\subset \C\setminus(-\infty,0]\), then \(Y\)
is linearly isomorphic to the complexification of a real Banach space.
\end{corollary}

\begin{proof}
Apply Theorem \ref{thm:spectral-normalization}.
\end{proof}

\begin{theorem}\label{thm:finite-bad-spectrum}
Let \(Y\) be a complex Banach space and let \(A:Y\to Y\) be a bounded
invertible anti-linear operator.  Put \(L=A^2\).  Suppose that
\[
  B=\sigma(L)\cap(-\infty,0]
\]
is finite and that every point of \(B\) is an isolated spectral value whose
Riesz spectral subspace is finite-dimensional.  Then \(Y\) is
complex-linearly isomorphic to the complexification of a real Banach space.
\end{theorem}

\begin{proof}
If \(B=\emptyset\), this is exactly Theorem
\ref{thm:spectral-normalization}.  Assume \(B\ne\emptyset\).

Let \(P_B\) be the Riesz spectral projection of \(L\) associated with \(B\).
Set \(E=P_BY\) and \(F=(I-P_B)Y\).  By hypothesis \(E\) is finite-dimensional,
and
\[
  Y=E\oplus F
\]
as a topological direct sum of complex subspaces, and the spectrum of \(L|_F\)
is contained in \(\C\setminus(-\infty,0]\).

The spaces \(E\) and \(F\) are invariant under \(A\).  Indeed, for every
\(z\in\C\),
\[
  A(zI-L)=(\overline z I-L)A .
\]
Choose the contours defining \(P_B\) symmetric with respect to complex
conjugation.  The Riesz integral and the anti-linearity of \(A\) then give
\(AP_B=P_BA\).  Therefore \(A|_F\) is a bounded anti-linear automorphism of
\(F\), and Theorem
\ref{thm:spectral-normalization} gives a bounded conjugation on \(F\).

The finite-dimensional space \(E\) also has a bounded conjugation, by choosing
any complex basis and conjugating coefficients.  Taking the direct sum of these
two conjugations gives a bounded conjugation on \(Y=E\oplus F\).  Lemma
\ref{lem:conjugation-real-form} completes the proof.
\end{proof}

\begin{corollary}\label{thm:compact-perturbation}
Let \(Y\) be a complex Banach space and let \(A:Y\to Y\) be a bounded
invertible anti-linear operator.  Suppose
\[
  A^2=\lambda I+K,
\]
where \(K\) is compact and \(\lambda\in\C\setminus(-\infty,0]\).  Then \(Y\)
is complex-linearly isomorphic to the complexification of a real Banach space.
\end{corollary}

\begin{proof}
Let \(L=A^2\).  The spectrum of the compact perturbation \(L=\lambda I+K\) is
the point \(\lambda\), together with at most countably many isolated eigenvalues
of finite algebraic multiplicity which can accumulate only at \(\lambda\).
Since \(\lambda\notin(-\infty,0]\), the bad spectral set
\(\sigma(L)\cap(-\infty,0]\) is finite, and all its Riesz spectral subspaces
are finite-dimensional.  Apply Theorem \ref{thm:finite-bad-spectrum}.
\end{proof}

\begin{corollary}\label{cor:scalar-compact-obstruction}
Let \(Y\) be a complex Banach space in which every bounded complex-linear
operator is of the form \(\mu I+K\), with \(K\) compact.  If \(Y\) is
anti-linearly isomorphic to itself and some anti-linear automorphism \(A\)
satisfies
\[
  A^2=\lambda I+K,\qquad \lambda\notin(-\infty,0],
\]
then \(Y\) has a real form.

Consequently, any non-real-form space of scalar-plus-compact type which is
anti-linearly isomorphic to itself must have only the remaining
``quaternionic'' obstruction: every such anti-linear square has scalar part on
the non-positive real axis.
\end{corollary}

\begin{proof}
The first assertion is Corollary \ref{thm:compact-perturbation}.  The final
sentence is just the contrapositive applied to anti-linear automorphisms of
\(Y\).
\end{proof}

\section{A Perturbative Projection Criterion}

\begin{theorem}\label{thm:near-real-projection}
Let \(Z\) be a complex Banach space with a bounded conjugation \(J\).  Let
\(P:Z\to Z\) be a bounded complex-linear projection and set \(Y=PZ\).  If
\[
  \|P\|\,\|P-JPJ\|<1,
\]
then \(Y\) is complex-linearly isomorphic to the complexification of a real
Banach space.
\end{theorem}

\begin{proof}
Put \(Q=JPJ\).  Then \(Q\) is a complex-linear projection onto \(JY\).  Define
the bounded anti-linear map
\[
  A:Y\to Y,\qquad Ay=PJy .
\]
Its square is the bounded linear operator
\[
  L=A^2=PQ|_Y .
\]
For \(y\in Y\), since \(Py=y\),
\[
  \|Ly-y\|=\|P(Q-P)y\|
    \leq \|P\|\,\|Q-P\|\,\|y\| < \|y\|.
\]
Thus \(\|L-I_Y\|<1\).  In particular \(L\) is invertible, and the binomial
series gives a bounded inverse square root
\[
  S=L^{-1/2}
\]
inside the closed algebra generated by \(L\).  Moreover \(A\) is injective
because \(A^2=L\) is invertible, and \(A\) is surjective since for any
\(y\in Y\),
\[
  y=A(A L^{-1}y).
\]
Thus \(A\) is an anti-linear automorphism.  The binomial coefficients are real,
so \(A\) commutes with \(S\).  Therefore
\[
  C=SA
\]
is bounded, anti-linear, and satisfies
\[
  C^2=SASA=S^2A^2=L^{-1}L=I_Y .
\]
So \(C\) is a bounded conjugation on \(Y\).  Lemma
\ref{lem:conjugation-real-form} again gives the desired real form.
\end{proof}

\begin{corollary}
Let \(Z=E_\C\) be a complex Banach lattice with its canonical conjugation
\(J\).  If \(P\) is a projection on \(Z\) whose range is \(Y\) and
\(\|P\|\,\|P-JPJ\|<1\), then \(Y\) is linearly isomorphic to the
complexification of a real Banach space.  In particular this holds whenever
\(P=JPJ\).
\end{corollary}

\section{Limitations And Novelty Check}

The result is only a partial positive answer to Remark 5.5.  It does not
settle arbitrary infinite-codimensional complemented subspaces of complex
Banach lattices, and it gives no counterexample to the full question.
Instead, it shows that any counterexample must avoid finite codimension, must
not have an unconditional Schauder basis, and must not admit either a
sufficiently conjugation-stable projection or an anti-linear self-isomorphism
whose square is spectrally normalizable, even up to compact perturbation away
from the non-positive real axis.

The run indexes were searched for arXiv:2310.02196, the title of the source
paper, and the phrases ``complexification'', ``complex conjugate'',
``GL-lust'', and ``complex Banach lattice''.  A web search on June 16, 2026
found the source arXiv page, the 2025 survey by de Hevia and Tradacete on
complemented subspaces of Banach lattices, Kalton's \(Z_2(\alpha)\) examples,
and Gordon--Junge--Meyer--Reisner's finite-dimensional analysis of the
Kalton--Peck twisted sum, but no exact prior run packet or later answer to
Remark 5.5.  The arguments above are elementary, so novelty confidence should
be treated as modest; the intended contribution is the scoped partial answer
and the explicit obstruction profile.

\begin{thebibliography}{9}

\bibitem{DHMST}
D. de Hevia, G. Martinez-Cervantes, A. Salguero-Alarcon, and P. Tradacete,
\emph{A counterexample to the complemented subspace problem in Banach
lattices}, arXiv:2310.02196v2, 2025.

\bibitem{HTsurvey}
D. de Hevia and P. Tradacete,
\emph{Complemented subspaces of Banach lattices}, arXiv:2505.24084v2, 2025.

\bibitem{Kalton}
N. J. Kalton,
\emph{An elementary example of a Banach space not isomorphic to its complex
conjugate}, arXiv:math/9402207, 1994.

\bibitem{GJMR}
Y. Gordon, M. Junge, M. Meyer, and S. Reisner,
\emph{The GL-l.u.st. constant and asymmetry of the Kalton-Peck twisted sum in
finite dimensions}, arXiv:1007.4692, 2010.

\end{thebibliography}

\end{document}
